\documentclass[../main.tex]{subfiles}

\begin{document}

\chapter{Tautologi, Kontradiksi, Ekivalensi Logis, dan Kontingensi}

\begin{subcpmk}
  \item Mahasiswa mampu mengidentifikasi serta mengklasifikasikan formasi nilai proposisi ke dalam kelas validitas matematis (tautologi, kontradiksi, atau kontingensi).
  \item Mahasiswa mampu membuktikan ekivalensi logis dari minimal dua pernyataan deklaratif program komputasi yang ekuivalen menggunakan manipulasi simbol.
\end{subcpmk}

\section{Klasifikasi Kategori Komputasi Logis}

Berdasarkan nilai mutlak pada rekapitulasi kolom paling kanan Tabel Kebesarannya (evaluasi akhirnya), semua proposisi yang majemuk \textit{(compound propositions)} dapat dengan tegar diklasifikasikan ke tiga strata konseptual yang eksklusif satu sama lain. 

\subsection{Tautologi}
Tautologi (bahasa Inggris: \textit{Tautology}) dapat disederhanakan konsepnya menjadi proposisi majemuk dengan fungsi evaluasi yang bernilai \textbf{MUTLAK BENAR (TRUE/T)} untuk \textit{segala/keseluruhan} kombinasi pemetaan nilai-nilai kebenaran sub-preposisi elementernya. Ibarat pernyataan universal yang tak terkalahkan tanpa peduli prasyarat sistem manapun.

Contoh ekspresi rasional yang paling mendasar: $p \lor \sim p$ (Apakah saat ini komputer sedang Nyala, ATAU komputer tidak Nyala). Kalimat logis ini pasti Tautologi absolut.

\subsection{Kontradiksi}
Lawan ontologis dari Tautologi adalah Kontradiksi. Ia dinistai selalu memberikan nilai \textbf{MUTLAK SALAH (FALSE/F)} pada tiap konstelasi \textit{truth assignment} untuk atom elemen independen di dalam sirkuit fungsi. \textit{Bug programming syntax rules semantik} yang umum adalah algoritma yang macet di *Logic Trap (Kontradiksi).*
Contoh logis matematis: $p \land \sim p$ (Hal itu tak mungkin dapat terjadi secara bersamaan, ini adalah paradoks / \textit{error handler}).

\subsection{Kontingensi}
Mayoritas masalah real komputasi bukanlah bersifat absolut, dan bernaung dalam kelas Kontingensi. Kontingensi disederhanakan deskripsinya di mana ia punya setidaknya 1 baris hasil T dan setidaknya 1 baris hasil evaluasi berstatus F pada hasil tabel fungsi terminal akhirnya. Kebenarannya "BERGANTUNG" dari keadaan alam semesta logis preposisi sub-inputnya (\textit{state dependent conditional output}).

\section{Ekivalensi Logis (Logical Equivalences)}
Pada pengoptimalan algoritma kode baris \textit{(Compilers Semantic Analysis)}, pengembang mencari penulisan blok kode yang tercepat dikomputasikan namun "bernilai keluaran akhir identik dan tak beda nilainya di memori komputer". Dua sintaks proposisi majemuk komputasi direkognisi saling \textbf{Ekivalen secara logis} sekiranya susunan kolom di dalam pemetaan evaluasi \textit{Truth Table} finalnya menunjukkan hasil pencerminan nilai (\textit{T/F}) identik dan setara per baris iterasi secara totalnya.
Disimbolisasi dengan bentuk ekuivalen tri-strip : $A \equiv B$.

\begin{contoh}
Hukum De Morgan: $\sim(p \lor q) \equiv \sim p \land \sim q$\\
Kedua formasi matematis ini dipastikan ekuivalen alias berdimensi ekivalensi logis, membantah disjungsi setara nilainya dengan pemutakhiran konjungsi dari elemen dinegasikan paralel.
\end{contoh}

\begin{aktivitas}
  \item Bentuklah argumen komputasi: $p \rightarrow q$. Apa pola ekivalensi logis untuk kalimat kondisional tersebut bila diubah ke dalam gabungan proposisi Disjungsi ($\lor$) atau Negasi dasar? (Petunjuk: Konversi Implikasi $\equiv \sim p \lor q$).
\end{aktivitas}

\begin{latihan}
  \item Analisis dan uji menggunakan deret validitas baris tabel kebenaran: "Buktikan secara algoritmik fomal bahwa $~ (p \rightarrow q) \equiv p \land \sim q$!"
\end{latihan}


\begin{checklist}
  \item Mengenali karakteristik kolom *output* absolut benar atau absolut palsu yang terasosiasi secara logik.
  \item Terbiasa memeriksa penyebutan Hukum De Morgan dalam perbaikan blok parameter syarat kondisi pemrograman \textit{(Condition Refactoring)}.
\end{checklist}

\end{document}
